\begin{definition}[Matrix product with a central input]
    \label{def:ldp_central_input}
    \lean{MPSPreparation.centralInputMatrix, MPSPreparation.mixedProductMap,
        MPSPreparation.reverseTransposeChain, MPSPreparation.mixedChain,
        MPSPreparation.mixedConfigurationEquiv,
        MPSPreparation.mixedConfigurationEquiv_apply}
    \leanok
    Choose a central site of a chain with $q=\ell+1+r$ sites. Write $A$ for
    its tensor and list the left and right subchains from their respective
    outer boundaries inward. Thus the right product $R(\sigma_R)$ is the
    transpose of the product in the original physical ordering. For a
    matrix $G$ on the virtual pairs, define
    \begin{align}
        F_{((\alpha,\beta),i,(\delta,\gamma)),x}
         & =A^i_{\beta\gamma}G_{(\alpha,\delta),x},
         & V_{(\sigma_L,i,\sigma_R),x}
         & =\sum_{\alpha,\beta,\gamma,\delta}
        L(\sigma_L)_{\alpha\beta}R(\sigma_R)_{\delta\gamma}
        F_{((\alpha,\beta),i,(\delta,\gamma)),x}.
        \label{eq:ldp_central_input}
    \end{align}
\end{definition}

\begin{theorem}[Coefficients of the central product]
    \label{thm:ldp_central_coefficients}
    \lean{MPSPreparation.mixedProductMap_apply,
        MPSPreparation.eval_reverseTransposeChain, MPSPreparation.eval_mixedChain}
    \leanok
    \uses{def:ldp_central_input}
    The coefficients in (\ref{eq:ldp_central_input}) satisfy
    \begin{align}
        V_{(\sigma_L,i,\sigma_R),x}
         & =\sum_{\alpha,\delta}
        \bigl(L(\sigma_L)A^iR(\sigma_R)^{\mathsf T}\bigr)_{\alpha\delta}
        G_{(\alpha,\delta),x}.
        \label{eq:ldp_central_coefficients}
    \end{align}
    Reversing the right site order and transposing each matrix gives
    precisely the original left-to-right physical matrix product.
\end{theorem}
\begin{proof}
    \leanok
    \uses{def:ldp_central_input}
    Expand the two matrix products in (\ref{eq:ldp_central_coefficients})
    and interchange the finite sums over $\beta,\gamma,\delta$.
    Transposition reverses the order of matrix multiplication.
\end{proof}

\begin{theorem}[Isometric factorization of virtual-pair coefficients]
    \label{thm:ldp_pair_coefficient_factorization}
    \lean{MPSPreparation.chainPairMatrix, MPSPreparation.sequentialEvaluation,
        MPSPreparation.sequentialEvaluation_isIsometry,
        MPSPreparation.exists_chainPairMatrix_factorization}
    \leanok
    For any subchain with positive bond dimension $D$, its coefficient
    matrix $T_{\sigma,(\alpha,\beta)}=(A_1^{\sigma_1}\cdots
        A_n^{\sigma_n})_{\alpha\beta}$ factors as $T=ES$.
    Here $E_{\sigma,a}=(Q_1^{\sigma_1}\cdots Q_n^{\sigma_n})_{1a}$,
    the outer bond has dimension one, every bond has dimension at most
    $D^2$, and each site map $Q_p$ is an isometry from its inward bond
    to its outward bond and physical site. In particular, $E^\dagger E=\Id$.
\end{theorem}
\begin{proof}
    \leanok
    \uses{thm:seqanc_sweep_matrix_remainder, thm:seqanc_support_propagation,
        thm:ldp_sequential_factorization_general}
    Apply successive decompositions to the initial row
    $\sum_\alpha\bra{\alpha,\alpha}$ and the matrices $\Id_D\otimes A_p^i$.
    Restrict the remainder to its nonzero rows. Preservation of inner
    products along the isometric sites gives $E^\dagger E=\Id$.
\end{proof}

\begin{definition}[Mixed sequential factorization]
    \label{def:ldp_mixed_sequential}
    \lean{MPSPreparation.HasMixedSequentialFactorization}
    \leanok
    Fix a central site of a chain of $q=\ell+1+r$ sites and a positive
    bond dimension $D$. A mixed sequential factorization of
    $V\colon\C^{D^2}\to(\C^d)^{\otimes q}$ consists of two families of local
    isometries $Q^L_p$ and $Q^R_p$, with outer bond dimensions
    $b^L_0=b^R_0=1$ and every bond dimension at most $D^2$, and an isometry
    $C\colon\C^{D^2}\to\C^{b^L_\ell}\otimes\C^d\otimes\C^{b^R_r}$.
    Define $(E_L)_{\sigma_L,a}=(Q^L_1(\sigma_{L,1})\cdots
        Q^L_\ell(\sigma_{L,\ell}))_{1a}$ and define $E_R$ in the same manner
    from the right boundary. The required contraction is
    \begin{align}
        V & =(E_L\otimes\Id_d\otimes E_R)C,
          & \sum_i(Q^L_p(i))^\dagger Q^L_p(i) & =\Id_{b^L_p},
          & \sum_i(Q^R_p(i))^\dagger Q^R_p(i) & =\Id_{b^R_p}.
        \label{eq:ldp_mixed_factorization}
    \end{align}
    The input of dimension $D^2$ is attached to the centre, rather than to
    either exterior boundary.
\end{definition}

\begin{theorem}[Two-sided factorization of a matrix product isometry]
    \label{thm:ldp_mixed_sequential_general}
    \lean{MPSPreparation.exists_mixed_isometric_factorization}
    \leanok
    \uses{def:ldp_central_input, def:ldp_mixed_sequential}
    If the matrix product $V$ of (\ref{eq:ldp_central_input}) is an isometry,
    it has a mixed sequential factorization about its chosen central site.
\end{theorem}
\begin{proof}
    \leanok
    \uses{thm:ldp_pair_coefficient_factorization, def:ldp_central_input}
    Independent inward decompositions of the left and right virtual-pair
    coefficient matrices give $L=E_L T_L$ and $R=E_R T_R$.
    Set $C=(T_L\otimes\Id_d\otimes T_R)F$.
    Matrix multiplication gives (\ref{eq:ldp_mixed_factorization}).
    Writing $E=E_L\otimes\Id_d\otimes E_R$, the identity
    $C^\dagger C=C^\dagger E^\dagger EC=V^\dagger V=\Id$
    proves that the central map is an isometry.
\end{proof}

\begin{theorem}[Mixed sequential factorization of a polar isometry]
    \label{thm:ldp_mixed_sequential_polar}
    \lean{MPSPreparation.mixedPolarIsoMatrix,
        MPSPreparation.mixedPolarIsoMatrix_eq_mixedProductMap,
        MPSPreparation.exists_mixed_sequential_polarIsoMatrix,
        MPSPreparation.exists_mixed_sequential_polarIsoMatrix_of_split}
    \leanok
    \uses{def:ldp_mixed_sequential, def:ldpolar_tensor, def:injective}
    Let $A_1,\ldots,A_q$ have positive bond dimension $D$, let $q\ge1$,
    and suppose the blocked tensor $B$ is injective. Its polar isometry
    $V=BP^{-1}$ has a mixed sequential factorization about any chosen
    physical site, including either endpoint. This is the injective case of the
    two-sided construction in \cite[footnote~4 to eqs.~(13)--(15)]{Malz2024Preparation}.
    For non-injective $B$, footnote~3 uses the pseudoinverse and
    $V^\dagger V$ is the polar support projector, rather than the identity;
    the corresponding mixed factorization on that support is not included here
    \cite{gap:mswc24_mixed_polar_injectivity_scope}.
\end{theorem}
\begin{proof}
    \leanok
    \uses{thm:ldp_mixed_sequential_general, thm:ldp_central_coefficients,
        thm:ldpolar_tensor_injective, lem:ldp_sequential_factorization_first_inputs}
    By Theorem~\ref{thm:ldpolar_tensor_injective}, $V^\dagger V=\Id$.
    Lemma~\ref{lem:ldp_sequential_factorization_first_inputs} expresses
    $V=BP^{-1}$ in the matrix-product coefficients of the given chain.
    Split the chain at the chosen site and orient its right part inward.
    Equation~(\ref{eq:ldp_central_coefficients}), with $G=P^{-1}$, gives
    these same coefficients. Apply the two-sided factorization theorem.
\end{proof}
